\section{Background and Related Work}

\subsection{Validation as a chemometric measurement problem}

Chemometric validation concerns the relation between the fitted model, the calibration sample, and the population to which predictions will be applied. Internal resampling is useful for estimating variability, but it does not by itself guarantee that the held-out observations represent a scientifically meaningful external prediction problem \citep{esbensen2010proper,tropsha2003earnest,gramatica2007principles}. Applicability-domain methods make the same point from the perspective of chemical coverage: prediction reliability depends on whether a test compound is represented by the training domain \citep{netzeva2005applicability,sahigara2012applicability}. More recent chemometric work has emphasized leakage, robustness, and the need to report validation design explicitly rather than treating a numerical score as self-interpreting \citep{lasfar2024robustness,kiraly2025leakage,ezenarro2025validation,camacho2026validation}.

For molecular machine learning, this implies that a split is part of the estimand. A model evaluated on chemically close random holdouts answers a different question from the same model evaluated on a structurally separated or temporally later test set. The relevant issue is therefore not which split is universally ``best'', but which properties of the test population are changed by the split and whether the resulting claim is stable to defensible protocol choices.

\subsection{Molecular benchmark splitting and out-of-distribution evaluation}

MoleculeNet established common molecular-property datasets and popularized random and scaffold-based evaluation \citep{wu2018moleculenet}. Scaffold partitions are commonly based on Bemis--Murcko frameworks \citep{bemis1996properties} and are intended to reduce direct framework overlap between training and test molecules. Comparative studies have nevertheless documented strong performance variability across split schemes and random seeds, as well as ambiguity in interpreting scaffold-held-out performance as generic chemical-space generalization \citep{yang2019learned,deng2023systematic}. Activity cliffs provide a complementary example in which local structural similarity does not imply similar response, further limiting what a single aggregate benchmark score can establish \citep{vantilborg2022cliffs,jimenez2022attribution}.

Several recent methods improve the realism or structural separation of molecular validation. SIMPD generates simulated time splits designed to resemble prospective medicinal-chemistry data \citep{landrum2023simpd}. DataSAIL formulates leakage-reduced splitting as a constrained optimization problem over similarity relationships and supports heterogeneous biological and molecular data \citep{joeres2025datasail}. Recent systematic QSAR work also shows that rational or similarity-based internal partitions can alter apparent performance without necessarily improving external generalization \citep{li2026partition}. These approaches address the important problem of defining a more appropriate deployment-like holdout.

Our question is different and complementary. We do not propose a new universal out-of-distribution splitter. Instead, we use a controlled pair of scaffold-disjoint partitions to isolate the contribution of one benchmark property---target-distribution mismatch---while holding the test size and search budget fixed. This turns split construction into a perturbation experiment rather than a leaderboard competition among unrelated partitioning rules.

\subsection{Target distribution as a benchmark property}

For classification, target-distribution shift appears as a change in class prevalence; for regression, it appears as a shift in the response distribution. Such changes can alter predictive difficulty even when structural separation is unchanged. Stratification is routinely used in ordinary classification resampling, but scaffold constraints make exact stratification nontrivial because complete structural groups must move together. A target-aware scaffold search can therefore improve class or response balance while simultaneously selecting a different region of chemical space. Without an exact-size target-blind comparator, the performance consequence of that intervention is confounded by other partition properties.

The present paired design separates the \emph{ability to balance the target} from the \emph{predictive consequence of doing so}. The target-balanced candidate is chosen from exactly the same target-blind candidate pool as its size-matched reference and must have the same number of test observations. The candidate-pool distribution is used to describe the geometry and selectivity of the search, not as an independent inferential null, because the balanced candidate is selected as an extremum from that pool.

\subsection{Molecular identity and scaffold semantics}

Chemical preprocessing can change what counts as a unique observation. Salts, counterions, mixtures, charge states, tautomeric forms, and duplicate measurements influence fingerprints, graph representations, deduplication, and structural grouping. Recent proposals for executable dataset contracts explicitly treat chemical identity rules, endpoint aggregation, split definitions, and leakage diagnostics as components of the QSAR claim \citep{nael2026dataset}. This view is particularly relevant to scaffold evaluation: preprocessing occurs before the scaffold identity is defined, so representation choices can propagate into the validation partition itself.

A second ambiguity occurs for acyclic molecules. The Bemis--Murcko framework of a molecule without a ring is empty. Implementations must therefore decide whether all acyclic compounds constitute one structural group, whether they are treated separately, or whether another grouping rule is used. These alternatives encode different meanings of ``unseen scaffold''. We explicitly test this semantic choice rather than assuming it is an implementation detail.

Together, the literature motivates a shift from architecture-centric benchmarking toward claim-centric validation. Our contribution is to make three usually hidden components measurable within one audit: target-distribution control under exact-size pairing, scaffold semantics for acyclic compounds, and disconnected-component molecular representation.
